\section{Experimental Evaluation}
We evaluate the proposed evidence-certification method along three axes: protocol behavior, decision semantics, and cryptographic cost. First, we manipulate the evidence, task context, and challenge parameters to determine whether the implementation handles invalid inputs as specified. We then control the true state, measurement errors, and Byzantine participation in synthetic workloads to assess the scope of the four-valued decision rule and compare it with alternative methods. Finally, we measure the proving costs of both Groth16 circuit backends and examine scalability alongside a public location-proof implementation and a resource-constrained compilation study. The tests evaluate distinct aspects of the protocol, finite model, and implementation; they do not constitute an end-to-end vehicle deployment.

\subsection{Implementation and experimental setup}
\subsubsection{Cryptographic implementation and test environment}
The prototype is implemented in Go 1.25.6 using gnark 0.14.0~\cite{gnark}, Groth16 over BN254~\cite{Groth2016}, and gnark's EdDSA/MiMC primitives for signing evidence and verifying signatures inside the circuit~\cite{MiMC2016}. The signature curve is \texttt{tedwards.BN254}, not Ed25519. Key generation, signing, parameter setup, proving, and verification all execute actual cryptographic implementations. BN254 was chosen for a reproducible prototype and benchmark; it is neither a post-quantum security claim nor a validated long-term deployment choice. The source and experiment scripts are maintained in the implementation repository. Its identifying link is omitted from this review copy; the source paths \texttt{prototype/} and \texttt{experiments/} refer to that repository's root.

Proof benchmarks ran on a Mac mini with an eight-core Apple M1 processor and 16 GB of memory. For each circuit configuration, one experimental setup was followed by 30 proofs and verifications using a fixed valid signed witness. Reported proving times exclude circuit compilation, evidence acquisition, network transport, queueing, and energy consumption. For context, we also implemented a simplified hidden-point region-membership circuit on the same gnark/Groth16/BN254 stack. Without issuer authentication, threshold fusion, or four-valued decisions, this circuit is a structural point of reference rather than a functionally equivalent baseline.

To examine circuit size separately from proving time, we compiled signed-evidence R1CS instances at issuer counts $n\in\{4,8,12,16\}$ and dimensions $d\in\{2,3\}$. For each pair we evaluated subset circuits with $k=n-1$, $k=n-2$, and $k=n/2$, along with the corresponding lower-corner circuit, for 32 compilation invocations in total. Each run had a 30-second timeout, a termination threshold of 3 GiB sampled resident set size (RSS), and a 2500 MiB Go soft memory target. We do not impute constraint counts to resource-limited runs. This experiment measures compilation feasibility and size, not proof latency at the largest configurations.

\subsubsection{Synthetic evidence and trajectory replay}
The semantic study uses the three-dimensional integer domain $\Omega=[0,63]^3$ and policy $P=[18,45]^3$. It configures five issuers, Byzantine budget $f=1$, and baseline honest-miscoverage allowance $r=0$. The design crosses maximum uncertainty radii $\{2,4,6\}$, actual Byzantine counts $\{0,1,2\}$, and three error models (independent calibrated, correlated calibrated, and undercovered), yielding 27 conditions. Each condition contains 30 random seeds and 80 events per seed, or 64,800 base observations in total. Three quarters of the events are placed near a policy face to emphasize boundary decisions. Honest interval radii are drawn from the ranges specified for each condition. Byzantine sources select their regions adversarially using the true state and public policy rather than only introducing random perturbations.

We compare the robust feasible-set rule with a median point claim, majority-center voting, a fixed maximum-radius margin, and all-source intersection. False acceptance (FA) uses truly noncompliant events as its denominator; false rejection (FR) uses truly compliant events. \emph{Decisive coverage} is the fraction of all events classified Valid or Invalid, excluding Ambiguous and Inconsistent decisions. It measures the frequency of definite answers, not their accuracy. When a seed has no events for a class-specific denominator, we do not assign an arbitrary zero error rate. Randomized equivalence tests also compare the native subset classifier, lower-corner search, and exhaustive classifier on tested inputs. The semantic workload uses the native classifier and does not generate a full Groth16 proof for every synthetic event; cryptographic timing is obtained from separate measurements.

Each event records $M$, the number of honest intervals that fail to cover the true state $z^\star$, allowing explicit checks of Proposition~\ref{prop:soundness}'s premise $M\le r$. To isolate threshold effects, the same 64,800 observations are re-evaluated with $r\in\{0,1,2\}$. This gives 194,400 \emph{paired evaluations}, not 194,400 new independent observations.

For trajectory replay, we extract 787 measured reference states spanning 786 seconds from UrbanNav-HK-Medium-Urban-1~\cite{UrbanNav}. After transformation to a local integer coordinate frame, a fixed percentile-defined policy includes 169 of the states. Each of 30 replay runs simulates four honest evidence sources, one policy-aware Byzantine source, and bounded correlated errors along the measured path. A stress mode multiplies the common-mode bias by four. The reference motion is measured; the multisource uncertainty boxes are simulated. This experiment therefore does not demonstrate that a real multisensor positioning system satisfies the physical coverage assumption.

\subsection{Protocol-security regression tests}
Table~\ref{tab:security-validation} records adversarial tests against the implemented proof-and-gateway path. Mutations target the evidence, decision code, context, issuer roster, and challenge nonce. The purpose is to check whether the implementation enforces the constraints from Sections~3 and~4, not to assess the physical accuracy of sensors.

\begin{table}[htbp]\centering\small
\caption{Adversarial regression tests of proof generation and gateway verification.}\label{tab:security-validation}
\begin{tabularx}{\linewidth}{@{}p{0.22\linewidth}p{0.35\linewidth}X@{}}\toprule
Attack & Enforced constraint & Observed prototype response\\\midrule
Evidence-bound tampering & Signed interval endpoints, issuer slot, context & Altered signed bounds were rejected\\
Decision substitution & Agreement between four-valued circuit output and public decision code & Tested mismatched proof/code pairs were rejected\\
Task or report substitution & Binding to task, report, threshold, roster, nonce and expiry & Tested cross-context substitutions were rejected\\
Challenge replay & Expiry and atomic state update after successful verification & Unknown or expired nonces were rejected; one of eight concurrent copies succeeded\\
Roster manipulation & Distinct ordered keys and trust-domain labels & Duplicated keys or labels were rejected\\
Missing issuer evidence & Complete evidence from the provisioned roster is required & Protocol aborted rather than certifying inconsistency\\\bottomrule
\end{tabularx}\end{table}

The results show that a proof associated with a different task or report cannot be transplanted directly into the current public statement, and missing roster data are not misreported as semantic inconsistency. The eight-way replay test supports only the implemented single-instance, no-restart guarantee of atomically marking a successfully used nonce. It does not establish distributed gateway safety, persistence across restarts, or business-report deduplication.
